\chapter{Methods}\label{chap:methods}

\section{Data sources and acquisition}\label{sec:data-sources}

Three data sources go into one hourly dataset for WTI crude oil: market data on prices and trading activity, structured macroeconomic indicators, and unstructured news. Each one arrives on its own clock and is accessed in its own way, so getting them onto a shared time grid is a real methodological problem rather than a formatting step. Section~\ref{sec:temporal-alignment} deals with it.

\subsection{Commodity selection}\label{sec:commodity-selection}

We settled on WTI crude oil futures after screening a few candidates. Sugar futures were the first alternative considered and were dropped: prices moved too little and reacted too weakly to news at hourly resolution, and the research questions here need both. WTI has the opposite profile. Its event structure is rich, driven by geopolitics, OPEC policy, the weekly EIA inventory cycle, and macroeconomic forces, and those events show up as hourly volume and price reactions we can actually measure.

\subsection{Market data}\label{sec:market-data}

Hourly OHLCV data are pulled with the yfinance Python library \citep{yfinance} for three instruments:

\begin{itemize}
  \item \texttt{CL=F}: WTI crude oil front-month futures contract, the primary target series
  \item \texttt{DX-Y.NYB}: the U.S.\ Dollar Index (DXY), used as a macroeconomic covariate
  \item \texttt{\^{}VIX}: the CBOE Volatility Index, used as a market-stress covariate
\end{itemize}

DXY and VIX come in as exogenous controls in the Phase~2 pipeline (Section~\ref{sec:phase2}), which lets us separate the news signal from wider macro-financial movement. They are not used in the Phase~1 analyses (Section~\ref{sec:phase1}), which look at news sentiment against trading volume directly, without macro controls.

From the raw OHLCV records we derive three liquidity and volatility proxies plus one return measure, all hourly:

\begin{itemize}
  \item \texttt{log\_volume}: the natural logarithm of contracts traded per hour, used as the primary liquidity measure
  \item \texttt{price\_range}: the hourly Parkinson range \citep{parkinson1980}, computed as $\log(\text{high}) - \log(\text{low})$, a high-frequency volatility proxy
  \item \texttt{log\_return}: the hourly log return, $\log(\text{close}_t / \text{close}_{t-1})$
  \item \texttt{amihud}: the Amihud illiquidity ratio at hourly resolution, computed as $\lvert \text{log\_return} \rvert / \text{volume}$ following \citet{amihud2002}. The ratio is undefined for hours with zero volume, so those rows are dropped from any analysis where Amihud is the dependent variable.
\end{itemize}

The canonical market grid runs from 13 May 2024 to 13 May 2026, which gives 11{,}232 hourly observations over roughly 24 months of trading. Phase~1 (Section~\ref{sec:phase1-data}) ran on an earlier pull that reached further back, roughly March 2024 to March 2026, for 11{,}219 aligned market hours, and we keep those values with the Phase~1 features rather than recomputing them. The reason the two windows differ is mundane: yfinance only serves hourly data for a rolling window of about the last two years. By the time we re-pulled the market data and extended the corpus through the post-war period for Phase~2 (Section~\ref{sec:phase2-data}), the earliest hours of 2024 had aged out of that window, so the Phase~2 grid starts in May 2024. The roughly 13-hour gap between the two counts comes from the re-pull and nothing else. No result depends on it and the methodology did not change.

\subsection{EIA inventory data}\label{sec:eia-data}

Weekly U.S.\ commercial crude oil inventory reports come from the U.S.\ Energy Information Administration (EIA) v2 API \citep{eia}. We query the WCRSTUS1 series (weekly ending stocks of crude oil and petroleum products, U.S.\ total) with \texttt{frequency=weekly}. We download the full history from January 2020 onward so that rolling baselines can be computed, even though only the 2024--2026 window is modelled.

Each weekly release becomes two derived features, available hourly once temporal alignment has been applied:

\begin{itemize}
  \item \texttt{is\_eia\_release}: a binary indicator that fires on the trading hour containing or immediately following the official release time
  \item \texttt{eia\_surprise}: the weekly inventory change relative to a rolling four-week baseline, capturing the directional surprise component of each release
\end{itemize}

EIA releases land on Wednesdays at 10:30 a.m.\ U.S.\ Eastern Time, and they are a scheduled event that the oil market watches closely, which is why we built features from them in the first place. In the reported models, though, the EIA cycle only enters through the \texttt{is\_wednesday} calendar proxy (Section~\ref{sec:input-typing}). We computed and stored \texttt{is\_eia\_release} and \texttt{eia\_surprise} and tested them in an extra model run, but the continuous surprise carried almost no feature importance and the release flag added little on top of the calendar proxy, so neither made it into the canonical model.

\subsection{News data}\label{sec:news-data}

News comes from the GDELT Project \citep{leetaru_schrodt2013}, a large machine-monitored corpus covering broadcast, print, and web sources worldwide. We use the GDELT 2.0 Document API to pull article metadata (URL, source, publication time, title) across a set of thematic queries aimed at WTI-relevant content.

The query set changed between the two phases. Phase~1 stuck to direct oil-market terms, giving eight queries:

\begin{lstlisting}
crude oil market
OPEC production cut
oil supply demand
petroleum inventory
Brent crude price
oil inventory report
energy market outlook
crude oil WTI price
\end{lstlisting}

The Phase~1 scrape returned 51{,}948 raw article records over the March 2024 to February 2026 window. Deduplication and English-only filtering cut that to 16{,}326 unique articles. To get body content we HTTP-fetched each URL and parsed the page, which worked about 80\% of the time. The failures clustered on paywalled domains, Cloudflare-protected sources, and pages that need JavaScript to render. An allow-list heuristic (Section~\ref{sec:filtering}) then judged 7{,}756 of the retrieved bodies to be substantive. Phase~1 analyses ran on a combined set of 13{,}690 articles: 7{,}756 with valid bodies going into title-plus-body sentiment extraction, plus 5{,}934 where the body fetch failed and we kept the title alone as a fallback.

For Phase~2 we widened the query set to catch the geopolitical and demand-side themes that dominated the post-war coverage period (Section~\ref{sec:phase2-data}):

\begin{lstlisting}
Iran sanctions oil
Saudi Arabia oil production
China oil demand
Russia oil exports
oil supply disruption
\end{lstlisting}

Re-running the scrape over the full window through May 2026 with both query sets returned 76{,}345 raw article records spanning January 2024 to May 2026. Deduplication and English-only filtering left 22{,}795 unique articles, which is the modeling-ready Phase~2 corpus (Section~\ref{sec:phase2-data}). Of those, 19{,}619 produced a real body. The other 3{,}176 came back as a null response, a Cloudflare or paywall block, or an explicit fetch error. Phase~2 pushes articles without bodies into LLM-based filtering instead of falling back to titles the way Phase~1 did, and Section~\ref{sec:filtering} explains why.

\subsection{Persistence}\label{sec:persistence}

Everything is stored in one SQLite database (\texttt{wti\_thesis.db}). The main tables are \texttt{articles} (raw GDELT records and scraped bodies), \texttt{market\_context} (the hourly grid of derived WTI liquidity features with DXY and VIX levels), \texttt{eia\_events} (raw EIA releases), \texttt{llm\_features} (LLM-extracted news features, Phase~2 only), \texttt{article\_entities} (canonical entity mentions per article, which is what the entity flags are built from), and \texttt{liquidity} (the modeling-ready joined view, populated by the alignment step in Section~\ref{sec:temporal-alignment}). The original Phase~1 pipeline passed CSV intermediates between notebooks, which let state drift out of sync. Putting everything in one relational database fixed that and made the canonical alignment a single deterministic join.

\section{Temporal alignment}\label{sec:temporal-alignment}

The three sources run on different clocks. Market data sits on a clean grid of trading hours, each bar tagged with its closing timestamp. EIA releases are weekly events at a fixed time. News articles are the awkward one: they can be published at any timestamp at all, including overnight, at weekends, and on holidays when nobody is trading. If we want a model that conditions liquidity at hour $t$ on news published before $t$, every article needs to be pinned to a specific hour on the market grid.

We handle this with two rules, both applied when articles are written to the database: ceiling assignment to the next trading hour, and forward-assignment of off-hours articles to the next available trading hour.

\subsection{Ceiling assignment to the next trading hour}\label{sec:ceiling}

For an article published at timestamp $\tau$, the assigned market hour is
\begin{equation}
  \text{datetime\_hour}(\tau) = \lceil \tau \rceil_{\text{1h}} .
\end{equation}
So an article published at 14:23 UTC goes to the 15:00 bar, not the 14:00 one. The reason is causal. The 14:00 bar covers 14:00 to 15:00, so part of its volume has already formed by 14:23. Assigning the article to the 14:00 bar would let news influence a candle that started forming before the news existed, which is backwards. The ceiling rule makes sure the article strictly precedes the start of the hour it is matched to.

An earlier version of the pipeline used \texttt{round} instead of \texttt{ceil}, which caused exactly this leak for any article published in the second half of an hour. We caught it early and the fix propagated through all the downstream notebooks.

The ceiling rule is conservative. It pushes news to the next full hour even when the article might well have moved the market within minutes, so it may understate the immediate intra-hour reaction. We accept that, because the alternative bias is worse: attributing pre-publication price movement to post-publication news would be a direct threat to the causal reading. The bias runs in the safe direction, understating rather than inflating the immediate reaction, and we come back to this where the lag results are interpreted (Section~\ref{sec:lag-ols}).

\subsection{Forward-assignment of off-hours articles}\label{sec:forward-assignment}

A lot of the corpus is published when no trading is happening: overnight, at weekends, or during holiday closures. There were two options. Drop those articles, or push them forward to the next available trading hour. We went with the second.

The reasoning is empirical. Geopolitical events, supply shocks, and weekend OPEC communiqu\'es get covered during off-hours as a matter of routine, and they are exactly the events that produce the biggest reactions once trading restarts. Dropping them would cut the most informative part of the corpus, and it would do so systematically rather than at random. Forward-assignment keeps that information. The price is a measurement-time lag between when an article was published and the hour it gets assigned to.

In practice, when an article's ceiling-assigned hour is not on the market grid because the market is shut, we assign it to the next available \texttt{market\_context} row in chronological order. This is a vectorised pandas \texttt{merge\_asof} join with \texttt{direction='forward'}, which does the whole corpus in one pass.

We keep a diagnostic column, \texttt{assignment\_gap}, which is the difference in hours between the assigned market hour and the original publication timestamp:
\begin{equation}
  \text{assignment\_gap} = \text{datetime\_hour} - \text{publication\_timestamp} .
\end{equation}
For an article published during a session this is at most one hour. For weekend or overnight articles it can run to many hours, or across a multi-day holiday closure in the worst case. The gap lives in the \texttt{liquidity} table and gets used downstream in two ways: as a control variable, and as the basis for stratified analyses separating \emph{contemporaneous} news (gap $<$ 2h) from \emph{forward-assigned} news (gap $\geq$ 2h). The exploratory Phase~1 VAR (later abandoned, Section~\ref{sec:var}) only used contemporaneous articles aggregated to hourly resolution. The Phase~2 TFT models use the whole corpus and let the model decide whether contemporaneous and forward-assigned articles deserve different weight.

\subsection{Joined output}\label{sec:joined-output}

Alignment produces the \texttt{liquidity} table, where each row is one article carrying the market-context features for its assigned hour. One boundary effect is worth flagging. Under forward-assignment (Section~\ref{sec:forward-assignment}), articles published before the market grid starts get mapped to the first available trading hour instead of being dropped. In the Phase~2 corpus that dumps roughly 3{,}600 pre-coverage articles (published before May 2024) onto the first grid hour, so the models see one unusually high-count hour right at the start of the series. Excluding pre-coverage articles outright gets rid of the artifact, and that adjustment is applied in the v2.3 ablation.

The resulting \texttt{liquidity} table has 22{,}795 rows: 15{,}290 contemporaneous (gap $<$ 2h) and 7{,}505 forward-assigned (gap $\geq$ 2h). Some analyses track the two separately, but the primary modeling combines them.

\section{News feature extraction approaches}\label{sec:feature-extraction}

News articles are text, and the models downstream need numbers. Getting from one to the other is the main NLP problem in this thesis. We do it two different ways across the two phases: FinBERT sentiment classification in Phase~1, and structured LLM extraction with Claude Haiku in Phase~2. They differ in what they output, and also in what they assume a ``news feature'' is in the first place.

\subsection{FinBERT sentiment extraction}\label{sec:finbert}

Phase~1 uses FinBERT \citep{araci2019}, a BERT variant fine-tuned on financial text for three-class sentiment (positive, neutral, negative). It is one of the most widely used sentiment models in financial NLP, which makes it a natural baseline.

We run it through the HuggingFace Transformers library \citep{wolf2020}, on Apple Silicon GPU acceleration via the MPS backend. For each article the model returns softmax probabilities over the three classes. The predicted label is the argmax, and we record a scalar \texttt{confidence} as the largest class probability:
\begin{equation}
  \text{confidence} = \max(P_{\text{positive}}, P_{\text{negative}}, P_{\text{neutral}}) .
\end{equation}

We keep the categorical label, all three class probabilities, and the confidence value for downstream use, persisted in the Phase~1 feature table \texttt{gdelt\_wti\_sentiment.csv}.

FinBERT runs twice per article:

\begin{itemize}
  \item \textbf{Title-only input}: the article title fed alone to the model
  \item \textbf{Title-plus-body input}: the title concatenated with the body, subject to FinBERT's 512-token input limit
\end{itemize}

Where the body has to be truncated, the title stays at the front of the sequence so the most salient part is never the thing that gets cut. Articles whose body retrieval failed (Section~\ref{sec:news-data}) fall back to title-only sentiment as their canonical input, and the body-availability status stays in the database so downstream analyses can tell the two regimes apart.

For the lag OLS analysis (Sections~\ref{sec:ols-spec} and~\ref{sec:phase1}) we use the FinBERT class probabilities directly rather than the argmax label, so each article is weighted by how confident the model was instead of being collapsed to a hard vote. The lag OLS takes \texttt{P(negative)} and \texttt{P(positive)} as two separate continuous regressors, which is what keeps the bearish-versus-bullish asymmetry that RQ2 needs. The earlier discrete mapping $\{\text{positive}: +1, \text{neutral}: 0, \text{negative}: -1\}$ throws the probability magnitudes away, and we keep it only as a comparison baseline in Section~\ref{sec:lag-ols}. The title-only probabilities are held back for the comparison in Section~\ref{sec:headline-bias}, the headline bias finding, where they support both a categorical and a continuous divergence measure.

\subsection{Structured LLM extraction with Claude Haiku}\label{sec:schema}

Phase~2 swaps FinBERT for Claude Haiku 4.5 \citep{anthropic_claude}, a smaller model in Anthropic's Claude family, and uses it to pull a richer set of structured features out of each article. Chapter~\ref{chap:experiments} (Section~\ref{sec:llm-extraction}) gives the full argument for the move. The short version is that FinBERT's three classes flatten something that is really continuous and multidimensional, and its confidence value (just the max class probability) does not carry much signal downstream.

The Haiku extraction went through one schema iteration. The first version (Haiku v1) produced the features for the first TFT training (Section~\ref{sec:tft-v1}) on the Phase~1 corpus. Per article it returned a continuous \texttt{sentiment\_score} in $[-1, +1]$ for the article's net directional impact on the WTI price ($-1$ strongly bearish, $+1$ strongly bullish, judged from the content rather than the headline tone), a \texttt{magnitude} score for how important the event is, a \texttt{price\_direction} label, an \texttt{event\_type} label, a list of \texttt{entities}, a \texttt{certainty} score, and a \texttt{time\_horizon} label. The second version (Haiku v2) came out of an inter-model calibration finding documented in Section~\ref{sec:calibration}. It dropped \texttt{price\_direction}, turned \texttt{event\_type} from a single label into an array of one to three labels ordered by salience, added an explicit \texttt{usable} boolean for filtering (Section~\ref{sec:filtering}), and added three orthogonal economic channel scores (\texttt{supply\_impact}, \texttt{demand\_impact}, \texttt{risk\_premium}), which Section~\ref{sec:channel-decomposition} describes. Appendix~\ref{app:prompt} has the full Haiku v2 prompt and tool schema.

\subsection{Tool-use API and schema enforcement}\label{sec:tool-use}

The Haiku extraction goes through Anthropic's tool-use API \citep{anthropic_tooluse}, with \texttt{tool\_choice} forcing the model to call a single \texttt{extract\_article\_features} tool. This is a methodological choice, not a stylistic one. Tool-use enforces the schema at the API level: enumerated values for categorical fields, type constraints on numeric fields, and required-field validation. An earlier pass that just asked the model in the prompt to produce valid JSON gave us schema violations constantly, for example returning \texttt{slightly\_bullish} where a numeric score belonged, or a \texttt{time\_horizon} label that was not on the list. Tool-use enforcement validates the numeric ranges and the scalar enumerated fields at the API boundary and gets rid of almost all of that. A small residue survives wherever the schema is permissive: the free-text \texttt{event\_type} array picked up an off-list label on 78 of 10{,}514 articles (0.74\%, \texttt{structural}), and one article has a malformed \texttt{time\_horizon} string. We treat both as noise downstream. So the enforcement removes the bulk of the parsing brittleness, not every last case.

The tool-use design also buys a short-circuit optimisation. When the model decides an article is unusable, whether that is a paywall block, a cookie notice, or off-topic content, only \texttt{usable} has to be set and every other field is optional. The model takes advantage of this and returns roughly 10 output tokens instead of 200 on unusable articles, which cuts the output-token cost of an unusable-article extraction by around 90\%.

\subsection{Channel decomposition}\label{sec:channel-decomposition}

Beyond the composite sentiment value, the Haiku schema has three continuous scores: \texttt{supply\_impact}, \texttt{demand\_impact}, and \texttt{risk\_premium}, each in $[-1, +1]$. The point of the three channels is to break the economic content of an article into dimensions that do not overlap:

\begin{itemize}
  \item \texttt{supply\_impact} measures the article's implication for the quantity of oil available to the market. Positive values indicate supply increases (sanctions lifted, OPEC adding barrels, new pipeline online); negative values indicate supply decreases (refinery outages, OPEC cuts, infrastructure attacks).
  \item \texttt{demand\_impact} measures the article's implication for oil consumption. Positive values indicate demand strength (Fed easing, China stimulus, growth surprise); negative values indicate demand weakness (recession signals, Fed hawkishness, lockdowns).
  \item \texttt{risk\_premium} measures the change in geopolitical or operational risk priced into oil. Positive values indicate escalation (military strikes on oil infrastructure, sanctions imposed); negative values indicate de-escalation (ceasefire, sanctions lifted).
\end{itemize}

The channels were not in the original Phase~2 schema. They were a response to a calibration failure on the composite \texttt{sentiment\_score}, where Haiku and a GPT-family reference model disagreed about the sign of sentiment on high-magnitude geopolitical events. Chapter~\ref{chap:experiments} documents the full reasoning: the calibration failure is in Section~\ref{sec:calibration}, and the decomposition with its re-calibration is in Section~\ref{sec:channel-response}. Methodologically the channels do two jobs here. They are higher-resolution features for TFT v2 (Section~\ref{sec:tft-v2}), and they give us something defensible to put in place of human-rater agreement when validating LLM-extracted features (Section~\ref{sec:calibration}).

We keep the composite \texttt{sentiment\_score} alongside the channels rather than replacing it. Three reasons. It keeps continuity with the Phase~1 headline bias experiment, which needs a single sentiment value. It lets us run ablations comparing the composite signal against the decomposed channels in the TFT (Section~\ref{sec:tft-v2}). And it hedges against the channels turning out to be individually noisy once run at full scale.

\subsection{Magnitude, event type, entities, certainty, and time horizon}\label{sec:other-fields}

The rest of the Haiku schema covers dimensions of news content that do not overlap with each other. \texttt{magnitude} is a continuous score in $[0, 1]$ for how market-moving an event is, with anchors given in the prompt (0.0 = negligible event, 1.0 = historic; most articles land in the 0.1--0.4 range, with a major OPEC cut anchored at 0.9 and a full-scale geopolitical crisis at 1.0). \texttt{event\_type} is an array of one to three labels from a fixed set of six (\texttt{geopolitical}, \texttt{supply}, \texttt{demand}, \texttt{macro}, \texttt{technical}, \texttt{other}), ordered by salience within the article. \texttt{entities} lists the named actors central to the article (countries, organisations, oil companies, key officials). \texttt{certainty} is a continuous score in $[0, 1]$ separating rumour, analyst forecast, and confirmed fact. \texttt{time\_horizon} is a label from \texttt{immediate} (hours to a day), \texttt{short\_term} (days to weeks), and \texttt{structural} (months or longer themes such as energy transition).

Every one of these encodes something FinBERT's three classes cannot, and each one enters TFT v2 either as a continuous input or as a categorical that the Variable Selection Network can weight on its own. The raw \texttt{entities} lists get normalised in post-processing onto a fixed canonical set of 71 actors using an alias mapping, so ``Tehran'' and ``Iranian'' both map to \texttt{Iran}. At hourly aggregation each hour gets a binary vector of length 71 marking which canonical entities the dominant article of that hour named, and that vector enters TFT v2 as 71 entity-flag features. Appendix~\ref{app:entities} has the full list. Section~\ref{sec:llm-extraction} discusses the rationale for each field, including why \texttt{event\_type} is a list rather than a single value and why entity normalisation happens in post-processing instead of in the prompt.

\section{Filtering strategies}\label{sec:filtering}

A good fraction of what GDELT returns is not substantive news. The usual offenders are paywall placeholders, cookie consent notices, Cloudflare blocks, scraping error pages, generic stock-pick listicles, and articles that matched a query keyword but have nothing to do with oil markets. These have to be filtered out. Left in, they would inject noise into the sentiment signal and bias the models downstream.

Each phase filters differently. Phase~1 used a regex/keyword heuristic on the article body. Phase~2 replaced it with an LLM-judged binary call on whether the content is usable, exposed as the \texttt{usable} field in the Haiku extraction. Both survive as columns in the database (\texttt{body\_valid} and \texttt{usable}), which lets us compare them directly in Section~\ref{sec:filter-comparison}.

\subsection{Regex/keyword heuristic (Phase 1)}\label{sec:regex-filter}

The Phase~1 filter, \texttt{body\_valid}, is an allow-list run over the article body. A body is valid only if it meets all four criteria:

\begin{itemize}
  \item Minimum length of 400 characters (excluding whitespace and HTML residues)
  \item Presence of at least one term from a domain-specific keyword list covering energy, financial, and geopolitical vocabulary
  \item Absence of press-release boilerplate (vendor blacklist for known PR-distribution patterns)
  \item An all-caps ratio below 30\%, computed over words of more than two characters in texts containing more than ten words (excluding scraped pages that are predominantly menu links, category headers, or stock-ticker listings)
\end{itemize}

The filter grew over the early part of the project. It started as a blacklist of known bad phrases (Cloudflare error text, cookie notices) and turned into the allow-list above. On the Phase~1 corpus of 16{,}326 articles with English text and a successful body fetch, it passed 7{,}756 articles (47.5\%). The rest were either kept with title-only sentiment as a fallback or dropped from the analysis, depending on the notebook. Section~\ref{sec:phase1-data} spells out which analyses used which subset.

The regex filter is fast, reproducible, and transparent, in the sense that anyone can read the rules and predict what they will do to a given article. Inspection also made its limits obvious. Some articles passed while being mostly boilerplate, long disclaimers or repeated company descriptions that happened to mention energy keywords. Others failed while carrying legitimate analyst commentary, either because the body came in just under the 400-character threshold or because the vocabulary did not line up closely enough with the energy keyword list.

That is what pushed us to a better filter in Phase~2.

\subsection{LLM-judged usable flag (Phase 2)}\label{sec:usable-filter}

The Phase~2 filter is the \texttt{usable} boolean in the Haiku schema, which asks the LLM itself whether the body is substantive WTI-relevant content. The prompt tells it to return \texttt{usable=false} for paywall placeholders, Cloudflare blocks, cookie notices, and anything clearly not about oil markets, macro, or oil-relevant geopolitics. When \texttt{usable=false}, the model is allowed and told to skip every other field, which is the short-circuit from Section~\ref{sec:tool-use}.

Three things motivated moving from \texttt{body\_valid} to \texttt{usable}. First, the LLM can read semantic content that a regex cannot match, so it can tell a real analyst piece from a press release that happens to be full of energy keywords. Second, the LLM's judgement runs on the same content model as the rest of the feature extraction, so the filter and the features agree with each other: an article marked usable will produce features that are themselves usable. Third, sending more articles to the LLM costs little under batch processing with prompt caching, so casting a wider net at the input does not create a resource problem.

\texttt{usable} becomes the primary content filter in Phase~2. There is also a stricter variant, \texttt{usable\_strict}, which additionally requires at least one of the three economic channels (\texttt{supply\_impact}, \texttt{demand\_impact}, \texttt{risk\_premium}) to be non-zero. That drops 1{,}161 articles which the model judged topical but gave no channel signal at all. The point of the stricter filter is not that those articles are mistakes, and Section~\ref{sec:filter-comparison} looks at what they actually are. It is that the channels are the Phase~2 feature contribution, so an article carrying no channel signal has nothing to contribute to them. \texttt{usable\_strict=1} is the canonical filter for TFT v2 training (Section~\ref{sec:filter-comparison}), and \texttt{usable=1} is what the filter comparison uses. \texttt{body\_valid} stays in the database for the comparison in Section~\ref{sec:filter-comparison} but is not applied as a filter in any modelling step.

\subsection{Comparative analysis}\label{sec:filter-comparison-methods}

Because both filter columns exist for the Phase~2 corpus, we can look at where they disagree. Of the 19{,}619 Phase~2 articles the LLM processed, 13{,}550 are flagged \texttt{body\_valid=1} and 11{,}675 are flagged \texttt{usable=1}. They disagree on a non-trivial share of articles in both directions: articles the regex rejected that the LLM accepted as substantive (the regex's false positives), and articles the regex accepted that the LLM rejected as boilerplate (its false negatives). Section~\ref{sec:filter-comparison} reports the full agreement analysis, with Cohen's kappa and worked examples of the disagreements.

\section{Statistical models}\label{sec:statistical-models}

The empirical work on news impact rests on lag-structured ordinary least squares (OLS) regressions, set out in Section~\ref{sec:ols-spec}. These are the canonical statistical evidence for both the lag question (RQ1) and the bearish-versus-bullish asymmetry question (RQ2). We also fitted an exploratory vector autoregression during development and abandoned it. The sentiment signal is sparse and event-driven, with more than half of all hours carrying no contemporaneous news at all, and the VAR could not identify significant dynamics on it. We do not develop it as a method here. Section~\ref{sec:var} notes what happened and how it pushed us toward the deep-learning model of Section~\ref{sec:tft-methods}.

\subsection{Lag OLS specification}\label{sec:ols-spec}

The lag OLS asks whether news sentiment published at hour $t$ predicts trading volume at hour $t + k$, for $k$ over a fixed set of lag horizons. At lag $k$ the specification is
\begin{equation}
  \text{log\_volume}_{t+k} = \beta_0 + \beta_1\, P(\text{negative})_t + \beta_2\, P(\text{positive})_t + \varepsilon_t ,
\end{equation}
where $P(\text{negative})_t$ and $P(\text{positive})_t$ are the FinBERT title-plus-body class probabilities for the article published at hour $t$. The neutral probability is the omitted reference, since the three probabilities sum to one and including all three would be perfectly collinear. So $\beta_1$ and $\beta_2$ give the expected difference in \texttt{log\_volume} at hour $t + k$ for a maximally confident bearish or bullish article (class probability of one) against a neutral one.

Using two class probabilities instead of one signed sentiment regressor is not a detail, it is what makes the asymmetry question answerable. A signed specification forces the bearish coefficient to be exactly the negative of the bullish one, which assumes a symmetric response before any estimation happens. The two-probability version lets them move independently, and that is what RQ2 is asking about. This continuous encoding is the primary Phase~1 specification. It generalises the earlier two-dummy encoding, which has the same two-coefficient structure but with hard 0/1 indicators, by weighting each article by FinBERT's confidence, so a 0.95-bearish article counts for more than a 0.51-bearish one. We keep the dummy version as a baseline in Section~\ref{sec:lag-ols}.

We estimate a separate regression at each lag $k \in \{1, 2, 3, 4, 6, 8, 12\}$. The set is dense at short horizons, where we expect information propagation to be detectable, and thinner at long ones, where we expect the effect to fade. At each lag we report the two coefficients, their standard errors, p-values, and the regression $R^2$. The $R^2$ will be small, and that is fine: news sentiment is one of many things driving hourly trading volume, and the goal is to find a structured pattern of coefficients across lags, not to explain as much variance as possible. What we treat as evidence of a real lag effect is a coefficient significant at the conventional $p < 0.05$ threshold, repeated across several adjacent lags, with a magnitude pattern that hangs together.

The contemporaneous case ($k = 0$) is reported separately as a baseline. At hour $t$ the article has only just been published, so an effect there means the news moved the market inside the same one-hour bar. It turns out to be small but statistically significant (Section~\ref{sec:lag-ols}). That is not a null result: it says sentiment carries detectable signal from the moment of publication, which is what makes it worth tracing across the later lags. Section~\ref{sec:lag-ols} reports it next to the lag regressions.

All regressions are estimated by ordinary least squares with the default standard error formula, and each observation is one article. The Phase~1 corpus of 13{,}690 articles yields somewhere between roughly 9{,}900 and 11{,}600 rows per lag regression, after dropping articles whose target hour $t + k$ falls outside the coverage window. Coverage drops as the lag gets longer. The regression treats these as a cross-section of articles rather than a sequential time series, so we do not apply heteroskedasticity- and autocorrelation-consistent (HAC) corrections. If articles published close together share unobserved common shocks, that assumption will understate the true uncertainty somewhat. The pattern we rely on should survive it: even with moderately wider confidence intervals, bearish coefficients peaking at lag $+6$ hours with bullish coefficients tracing the same shape at lower magnitude would still be there.

\section{Temporal Fusion Transformer}\label{sec:tft-methods}

The Temporal Fusion Transformer (TFT) \citep{lim2021}, from Lim, Ar\i k, Loeff, and Pfister, is a deep-learning architecture for multi-horizon time-series forecasting with a mix of continuous and categorical inputs. It is our primary deep-learning model for Phase~2. We train two instances, TFT v1 and TFT v2, which share an architectural skeleton but differ in their features and their training data. This section covers the architecture. The feature configurations are in Sections~\ref{sec:tft-v1} and~\ref{sec:tft-v2}.

\subsection{Architecture overview}\label{sec:tft-architecture}

The TFT puts four components together into one end-to-end model:

\begin{itemize}
  \item A \textbf{Variable Selection Network (VSN)} at each input time step, which assigns a learned scalar weight to each input feature, allowing the model to suppress uninformative features automatically
  \item A pair of \textbf{LSTM encoder--decoder layers} that process the history (encoder) and the forecast horizon (decoder) at fixed resolution
  \item An \textbf{interpretable multi-head self-attention} layer over the encoded sequence, exposing which historical time steps the model attends to when producing each forecast
  \item A \textbf{quantile output head} that produces multiple quantile predictions in a single forward pass, enabling prediction intervals without separate models
\end{itemize}

Three of these give us something interpretability-wise. The Variable Selection Network produces per-feature importance scores that can be compared across features, which answers the question of which inputs the model is actually using. The attention layer produces per-time-step weights that lay out the temporal structure of the model's reasoning, and that goes straight to RQ1 on the lag structure. The quantile output gives us uncertainty estimates without training separate models for them.

\subsection{Input typing}\label{sec:input-typing}

The TFT sorts inputs into four categories:

\begin{itemize}
  \item \textbf{Time-varying known reals}: features whose values are known at every time step in both the history and the forecast horizon. Ours are calendar features (\texttt{hour}, \texttt{day\_of\_week}, \texttt{month}) and binary regime indicators (\texttt{is\_us\_session}, \texttt{is\_wednesday} as a proxy for the EIA release day).
  \item \textbf{Time-varying unknown reals}: features known only over the history, which have to be predicted for the forecast horizon. These are the target itself (\texttt{log\_volume}), the market-context features (\texttt{price\_range}, \texttt{log\_return}, \texttt{amihud}, \texttt{dxy}, \texttt{vix}), the continuous Haiku features (the composite \texttt{sentiment\_score}, the three channels, \texttt{magnitude}, \texttt{certainty}), and, in TFT v2, the 71 binary entity flags. Sections~\ref{sec:tft-v1} and~\ref{sec:tft-v2} give the per-instance feature lists.
  \item \textbf{Time-varying unknown categoricals}: categorical news features known only over the history. In TFT v2 these are the dominant article's \texttt{event\_type} and \texttt{time\_horizon}, passed as learned embeddings rather than integer codes. Whether they are encoded as proper categoricals or as integers is a design axis, and the v2 ablation isolates its effect (Section~\ref{sec:feature-importance}).
  \item \textbf{Static categoricals}: features fixed for the whole series. With a single asset this comes down to one asset identifier (\texttt{asset = 'WTI'}).
\end{itemize}

The Variable Selection Network learns separate gating weights per category, so a feature that is informative in the history as an unknown real but useless as a calendar input gets weighted correctly in both contexts.

\subsection{Forecast configuration}\label{sec:forecast-config}

Both instances condition each forecast on a fixed 48-hour encoder window, which is two full daily cycles of context. The length is deliberate. The Phase~1 lag OLS (Section~\ref{sec:lag-ols}) put the news-impact peak at lag $+6$ hours, and 48 hours comfortably covers that peak plus the second-day ``memory'' effects we saw in earlier TFT runs.

Target and horizon differ between the instances. TFT v1 (Section~\ref{sec:tft-v1}) predicts one target, \texttt{log\_volume}, one hour ahead. TFT v2 (Section~\ref{sec:tft-v2}) predicts three targets jointly, \texttt{log\_volume}, \texttt{amihud}, and \texttt{price\_range}, at four horizons (1, 3, 6, and 12 hours), so its prediction length is 12 hours and it emits a whole trajectory rather than one step. The $+6$ hour horizon is there to line up with the lag OLS peak.

Both are trained with a quantile loss, which gives point predictions at several quantiles in one forward pass and handles target outliers without needing distributional assumptions. TFT v1 uses a single \texttt{QuantileLoss} on its one target. TFT v2 uses a \texttt{MultiLoss} summing a \texttt{QuantileLoss} over each of the three targets. The median quantile is the point forecast throughout.

\subsection{Architectural hyperparameters}\label{sec:tft-hyperparams}

The TFT has a small set of hyperparameters controlling capacity. We hold the core width parameters constant across both instances, so that any difference between them comes from features and forecast configuration rather than from raw capacity:

\begin{itemize}
  \item \texttt{hidden\_size = 32}: the internal representation dimension
  \item \texttt{attention\_head\_size = 4}: the number of parallel attention heads
  \item \texttt{hidden\_continuous\_size = 16}: the embedding dimension for continuous variable processing
  \item \texttt{learning\_rate = 1e-3}: the initial Adam optimizer learning rate, with on-plateau reduction (patience 3)
\end{itemize}

Two things do differ. \texttt{dropout} goes from \texttt{0.1} in v1 to \texttt{0.15} in v2, a small bump in regularization for v2's bigger input space. And even with the width unchanged, the trainable-parameter count goes from roughly 113{,}000 in v1 to 298{,}329 in v2, because v2 brings in the 71 entity-flag inputs, two more targets, and the multi-horizon output head. Keeping the width fixed while the feature set and forecast configuration grow is the point: it keeps the instances comparable on architectural capacity, so what we are measuring is the effect of the Phase~2 feature work.

\subsection{Training procedure}\label{sec:training}

Training runs in PyTorch Lightning \citep{pytorch_lightning} through the \texttt{pytorch-forecasting} library \citep{pytorch_forecasting}. Both instances use the same optimizer and regularization scaffolding: Adam at the learning rate above, gradient clipping at 0.1 to keep updates stable over long sequences, early stopping on validation loss with a minimum improvement of \texttt{1e-4}, and a Google Colab T4 GPU. The \texttt{pytorch-forecasting} \texttt{TimeSeriesDataSet} wrapper builds the fixed-length sliding-window samples, encodes categoricals consistently, constructs the input tensors, and applies group normalization. Held-out sets go through the same wrapper with \texttt{stop\_randomization=True} so evaluation is deterministic.

Three training settings differ, reflecting the difference in scale:

\begin{itemize}
  \item \textbf{Temporal split.} TFT v1 uses an 80/20 split, the earlier 80\% of hours for training and the most recent 20\% for validation, with no held-out test set. TFT v2 uses a 60/20/20 split into train, validation, and test, with 48-hour buffers between partitions so encoder windows cannot leak across boundaries. The test set is what makes the pre-war versus war-regime evaluation of Section~\ref{sec:predictive-performance} possible.
  \item \textbf{Early stopping and epoch cap.} TFT v1 uses a patience of 5 epochs and a maximum of 30. TFT v2 raises these to 10 and 75. In practice early stopping halts both well before the cap.
  \item \textbf{Target normalization.} TFT v1 normalizes its single target with a \texttt{GroupNormalizer}. TFT v2 is multi-target, so it wraps one \texttt{GroupNormalizer} per target in a \texttt{MultiNormalizer}, scaling each of the three targets independently before the per-target losses are combined.
\end{itemize}

For reproducibility, the reported TFT v2 run also fixes the random seed (42) and uses deterministic algorithms.

\subsection{Evaluation and interpretation outputs}\label{sec:evaluation-methods}

Every trained TFT gives us a set of evaluation and interpretability outputs, all reported in Chapter~\ref{chap:experiments}.

\textbf{Performance against a persistence baseline.} The headline number is not the raw validation loss. The two instances train with different loss functions on different splits and horizons, so their loss values cannot be compared directly (Section~\ref{sec:phase-comparison}). Instead we score forecasts by mean absolute error (MAE) against a persistence baseline, the naive forecast that carries the current value forward to the target horizon ($\hat{y}_{t+k} = y_t$), which is a standard reference for financial time series. What we report is the percentage reduction in MAE over persistence, which normalizes for how hard each target and horizon is. For TFT v2 we compute these for every target (\texttt{log\_volume}, \texttt{amihud}, \texttt{price\_range}) and horizon (1, 3, 6, 12 hours) on four slices: the validation set, the full test set, and the pre-war and war portions of the test set, split at the Iran war onset (Section~\ref{sec:phase2-data}). The pre-war and war slices are how we measure whether a model trained pre-war generalizes across the regime boundary.

\textbf{Interpretability.} Two outputs are read straight off each model, both averaged over the evaluation set: the Variable Selection Network feature-importance ranking, which shows which inputs the model leaned on, and the multi-head attention weights, which give a per-lag map of which historical hours it attended to.

\textbf{Research-question diagnostics.} The attention output matters most for RQ1. If TFT attention peaks at a lag consistent with the Phase~1 lag OLS result, then two analyses that start from completely different places agree: the OLS finds a structured lag effect from a series of single-lag linear specifications, while the TFT learns one nonlinear function over the whole history and exposes the same structure through attention. For RQ2 we run a directional asymmetry test on the forecasts. Hours are split into bearish and bullish groups by the sign of their sentiment score, and a Welch t-test compares mean predicted \texttt{log\_volume} between the groups at each horizon and slice. This is deliberately a different estimand from the lag OLS coefficient comparison: it contrasts the average level of predicted volume across two groups of hours, where the OLS measures the marginal sensitivity of volume to sentiment. Section~\ref{sec:asymmetry} reads the two together with that distinction kept in view.

\section{Inter-model calibration methodology}\label{sec:calibration-methods}

Everything in Phase~2 is built on the Haiku-extracted features of Section~\ref{sec:feature-extraction}. Before leaning on them we want an external check that the extraction is sensible and that the LLM is not systematically wrong on the cases that matter most. This section is the methodology for that check. The results, and the schema change they prompted, are in Chapter~\ref{chap:experiments} (Sections~\ref{sec:calibration} and~\ref{sec:tft-v1}).

\subsection{The validation problem and our approach}\label{sec:validation-problem}

Normally you validate an NLP extraction against human-annotated ground truth on a labelled subset. That was not available here. Annotating WTI news properly takes domain expertise in oil markets and trading that the author does not have, and one non-expert annotator working at scale would add noise of unknown shape to the validation set rather than remove it. Hiring an expert at scale was beyond the project's resources.

So we use inter-model calibration instead, leaning on the evidence that LLMs are competent annotators in place of human labels \citep{gilardi2023} and can judge other models' output \citep{zheng2023}. The same articles, with the same prompt and tool schema, are scored by Claude Haiku 4.5 and by a GPT-family reference model (OpenAI's GPT-5.5, through the ChatGPT interface). Disagreement between two models from different developer families is not the same signal as disagreement between humans, but it is still a signal. The argument goes as follows.

Each model carries its own implicit theory of how to read news, shaped by its training data, its alignment procedure, and the inductive biases of its developer. These two share little in either training data composition or alignment style. When they agree, that agreement is about content robust to two independently trained readings of financial news. When they disagree, it means either the article is genuinely ambiguous or one of the models is systematically biased.

This is weaker than expert-annotated ground truth and we are not claiming otherwise: two models agreeing on a wrong reading would sail straight through the comparison. Chapter~\ref{chap:discussion} takes the limitation up as part of the discussion. It is still stronger than no external validation at all, and picking models from different families gets us as much diversity as possible given that both annotators are LLMs.

\subsection{Sample design}\label{sec:calibration-sample}

The calibration sample is drawn by stratified random sampling from the full Phase~2 corpus. The strata come from the Phase~1 regex body-validity classification, in three buckets:

\begin{itemize}
  \item \textbf{valid}: articles passing the regex \texttt{body\_valid=1} filter (Section~\ref{sec:regex-filter})
  \item \textbf{invalid\_substantial}: articles failing \texttt{body\_valid=1} but whose body field still contains substantial text content (length above a minimum threshold)
  \item \textbf{invalid\_weak}: articles failing \texttt{body\_valid=1} with bodies short enough to be likely scraping artefacts
\end{itemize}

We sample ten articles uniformly at random from each stratum, for a calibration sample of thirty. Stratifying this way makes sure the calibration sees the full range of content the LLM meets in production, including the marginal cases where the regex filter and the LLM are most likely to part ways.

\subsection{Extraction procedure}\label{sec:calibration-procedure}

Both models get the same system prompt and tool schema (Appendix~\ref{app:prompt}) for each of the thirty articles. For Claude Haiku this is the production extraction code from Section~\ref{sec:schema}. For the GPT reference, the prompt and schema go through the ChatGPT interface and the JSON output is extracted by hand.

The two output sets are then lined up by article identifier and compared field by field. The agreement metrics are computed symmetrically, so no metric depends on which model we call the reference.

\subsection{Agreement metrics}\label{sec:calibration-metrics}

For each field we compute one or more agreement measures, picked to suit the field's type:

\begin{itemize}
  \item For binary fields (\texttt{usable}): agreement rate as a count out of thirty, and Cohen's kappa to correct for chance agreement
  \item For continuous fields (\texttt{sentiment\_score}, \texttt{supply\_impact}, \texttt{demand\_impact}, \texttt{risk\_premium}, \texttt{magnitude}, \texttt{certainty}): Pearson correlation across all pairs, mean absolute difference (MAD), and the count of \textbf{sign disagreements}, meaning cases where the two models' scores have opposite signs (one positive, one negative). Near-zero scores within a small range around zero count as agreement on neutrality and do not add to the sign-disagreement count.
  \item For categorical fields (\texttt{event\_type}, \texttt{time\_horizon}): exact-match agreement rate, and for the array-valued \texttt{event\_type}, the proportion of cases where the top-ranked category from one model appears anywhere in the other model's list
\end{itemize}

Of these, the sign-disagreement count matters most. A small mean absolute difference can hide cases where the two models disagree about the qualitative direction of an effect, one calling an event bullish and the other bearish, and that is exactly the disagreement that would flow into downstream models as noise. The sign-disagreement count puts those cases on the surface.

\subsection{Use of the calibration outcome}\label{sec:calibration-use}

The calibration sample is small ($n = 30$) on purpose. A bigger sample means proportionally more annotation cost on the GPT side, and what we want is to spot qualitative patterns of agreement and disagreement, not to estimate a population agreement rate precisely. Thirty articles are enough to surface systematic disagreements, show what kind they are, and tell us whether to accept the schema or revise it.

If the calibration turns up significant disagreement on a field, the response is to work out where the disagreement comes from, revise the prompt or schema, then re-run the calibration on the same thirty articles with the new setup. Chapter~\ref{chap:experiments} documents this loop: Section~\ref{sec:calibration} has the first calibration result, and Section~\ref{sec:channel-response} has the schema revision and the re-run. Reusing the same thirty articles for both rounds is intentional, since it separates the effect of the schema change from the effect of sampling different articles.
